% ============================================================
\chapter*{Pr\'eface}
\addcontentsline{toc}{chapter}{Pr\'eface}
% ============================================================

Les \'equations aux d\'eriv\'ees partielles (EDP) constituent l'un des piliers
fondamentaux des math\'ematiques appliqu\'ees et de la physique math\'ematique.
Depuis les travaux fondateurs d'Euler, d'Alembert, Fourier et Laplace au
XVIII\ieme{} si\`ecle, elles n'ont cess\'e de jouer un r\^ole central dans la
mod\'elisation des ph\'enom\`enes naturels~: diffusion de la chaleur, propagation
des ondes, \'ecoulement des fluides, \'elasticit\'e des mat\'eriaux, \'electromagn\'etisme,
et plus r\'ecemment en finance, biologie et traitement d'images.

\section*{Objectifs du cours}

Ce cours s'adresse aux \'etudiants de troisi\`eme ann\'ee de licence (L3) en
math\'ematiques. Il suppose acquis les pr\'erequis suivants~:
\begin{itemize}
    \item Analyse r\'eelle~: int\'egration de Lebesgue, espaces $L^p$, convergences.
    \item Analyse fonctionnelle \'el\'ementaire~: espaces de Banach, espaces de Hilbert.
    \item \'Equations diff\'erentielles ordinaires (EDO)~: th\'eor\`eme de Cauchy-Lipschitz.
    \item Alg\`ebre lin\'eaire~: diagonalisation, formes bilin\'eaires.
\end{itemize}

L'objectif principal est de fournir une introduction rigoureuse mais accessible aux
m\'ethodes classiques et modernes de r\'esolution des EDP. Plus pr\'ecis\'ement, \`a
l'issue de ce cours, l'\'etudiant sera capable de~:
\begin{enumerate}
    \item Classifier les EDP du second ordre (elliptiques, paraboliques, hyperboliques)
          et comprendre les implications physiques de chaque type.
    \item R\'esoudre explicitement des EDP classiques par la m\'ethode des
          caract\'eristiques, la s\'eparation des variables et les s\'eries de Fourier.
    \item \'Enoncer et d\'emontrer les propri\'et\'es qualitatives fondamentales~:
          principe du maximum, r\'egularit\'e, comportement asymptotique.
    \item Utiliser la transform\'ee de Fourier pour r\'esoudre des EDP sur $\R^n$.
    \item Comprendre le cadre fonctionnel moderne~: espaces de Sobolev et
          formulations variationnelles.
\end{enumerate}

\section*{Organisation du cours}

Le cours est organis\'e en dix chapitres, suivant une progression p\'edagogique
naturelle~:

\begin{description}
    \item[Chapitre 1~:] \textbf{Introduction et classification.} On y pr\'esente
        le vocabulaire de base (ordre, lin\'earit\'e, conditions aux limites) et la
        classification des EDP du second ordre.

    \item[Chapitre 2~:] \textbf{\'Equation de transport.} Premi\`ere EDP \'etudi\'ee
        en d\'etail, elle introduit la m\'ethode des caract\'eristiques dans un
        cadre simple.

    \item[Chapitre 3~:] \textbf{S\'eparation des variables et s\'eries de Fourier.}
        On d\'eveloppe l'outil central de la s\'eparation des variables, appuy\'e
        sur la th\'eorie des s\'eries de Fourier.

    \item[Chapitre 4~:] \textbf{\'Equation de la chaleur.} \'Etude compl\`ete~:
        solution fondamentale, principe du maximum, r\'egularit\'e,
        comportement asymptotique.

    \item[Chapitre 5~:] \textbf{\'Equation des ondes.} Formule de d'Alembert,
        principe de Huygens, \'energie et vitesse finie de propagation.

    \item[Chapitre 6~:] \textbf{\'Equation de Laplace et fonctions harmoniques.}
        Propri\'et\'es de valeur moyenne, r\'egularit\'e, in\'egalit\'e de Harnack.

    \item[Chapitre 7~:] \textbf{Principe du maximum.} Versions faible et forte
        pour les op\'erateurs elliptiques et paraboliques.

    \item[Chapitre 8~:] \textbf{Transform\'ee de Fourier et applications.}
        La transform\'ee de Fourier comme outil de r\'esolution sur $\R^n$.

    \item[Chapitre 9~:] \textbf{Introduction aux espaces de Sobolev.}
        D\'erivation faible, espaces $H^1$ et $H^1_0$, in\'egalit\'es de Sobolev
        et de Poincar\'e.

    \item[Chapitre 10~:] \textbf{Formulation variationnelle.}
        Th\'eor\`eme de Lax-Milgram, probl\`emes elliptiques, introduction aux
        \'el\'ements finis.
\end{description}

\section*{Conventions et notations}

Tout au long de ce cours, nous utilisons les notations suivantes~:
\begin{itemize}
    \item $\Omega$ d\'esigne un ouvert born\'e de $\R^n$ de fronti\`ere $\partial\Omega$
          suffisamment r\'eguli\`ere (au moins $C^1$ par morceaux, sauf mention contraire).
    \item $u_t = \frac{\partial u}{\partial t}$, $u_x = \frac{\partial u}{\partial x}$,
          $u_{xx} = \frac{\partial^2 u}{\partial x^2}$.
    \item $\Delta u = \sum_{i=1}^{n} \frac{\partial^2 u}{\partial x_i^2}$ est le
          laplacien de $u$.
    \item $\nabla u = \left(\frac{\partial u}{\partial x_1}, \ldots,
          \frac{\partial u}{\partial x_n}\right)$ est le gradient.
    \item $Du$ d\'esigne la matrice jacobienne, $D^2 u$ la matrice hessienne.
    \item $C^k(\Omega)$, $C^\infty(\Omega)$, $C_c^\infty(\Omega)$~: espaces de
          fonctions de classe $C^k$, lisses, lisses \`a support compact.
    \item $L^p(\Omega)$~: espace de Lebesgue, $1 \leq p \leq \infty$.
    \item $H^s(\Omega) = W^{s,2}(\Omega)$~: espace de Sobolev.
    \item $\inner{\cdot}{\cdot}$ d\'esigne le produit scalaire dans $L^2$ ou dans un
          espace de Hilbert.
    \item $\norm{\cdot}$ d\'esigne une norme (le contexte pr\'ecise laquelle).
    \item Multi-indices~: pour $\alpha = (\alpha_1, \ldots, \alpha_n) \in \N^n$,
          $\abs{\alpha} = \alpha_1 + \cdots + \alpha_n$ et
          $D^\alpha u = \frac{\partial^{\abs{\alpha}} u}{\partial x_1^{\alpha_1}
          \cdots \partial x_n^{\alpha_n}}$.
\end{itemize}

\section*{Comment utiliser ce cours}

Chaque chapitre contient~:
\begin{itemize}
    \item Des \textbf{d\'efinitions} et \textbf{th\'eor\`emes} avec d\'emonstrations
          d\'etaill\'ees.
    \item Des \textbf{exemples} trait\'es compl\`etement pour illustrer les m\'ethodes.
    \item Des encadr\'es \og \textbf{Intuition} \fg pour d\'evelopper la compr\'ehension
          g\'eom\'etrique et physique.
    \item Des encadr\'es \og \textbf{Attention} \fg pour signaler les pi\`eges courants.
    \item Des \textbf{exercices} de difficult\'e croissante, certains avec indications.
\end{itemize}

Les d\'emonstrations marqu\'ees d'un ast\'erisque (*) sont hors programme et
propos\'ees pour les \'etudiants souhaitant approfondir.

\section*{Motivation physique}

Les EDP apparaissent naturellement d\`es que l'on mod\'elise un ph\'enom\`ene
d\'ependant de plusieurs variables ind\'ependantes. Consid\'erons quelques exemples
fondamentaux~:

\paragraph{Diffusion de la chaleur.}
Si $u(x,t)$ repr\'esente la temp\'erature en un point $x$ \`a l'instant $t$ dans
un corps conducteur, la loi de Fourier affirme que le flux de chaleur est
proportionnel au gradient de temp\'erature~: $\vec{q} = -k\nabla u$. Combin\'ee
avec la conservation de l'\'energie, on obtient l'\'equation de la chaleur~:
\[
    \frac{\partial u}{\partial t} = k \Delta u.
\]

\paragraph{Propagation des ondes.}
La vibration d'une corde, d'une membrane ou les ondes acoustiques sont d\'ecrites
par l'\'equation des ondes~:
\[
    \frac{\partial^2 u}{\partial t^2} = c^2 \Delta u,
\]
o\`u $c$ est la vitesse de propagation.

\paragraph{\'Electrostatique.}
Le potentiel \'electrique $u$ dans une r\'egion sans charge satisfait l'\'equation
de Laplace~:
\[
    \Delta u = 0.
\]
En pr\'esence d'une densit\'e de charge $\rho$, on obtient l'\'equation de Poisson
$\Delta u = -\rho/\varepsilon_0$.

\paragraph{M\'ecanique des fluides.}
Les \'equations de Navier-Stokes, qui d\'ecrivent le mouvement d'un fluide
visqueux incompressible, forment un syst\`eme d'EDP non lin\'eaires~:
\[
    \frac{\partial \vec{v}}{\partial t} + (\vec{v} \cdot \nabla)\vec{v}
    = -\frac{1}{\rho}\nabla p + \nu \Delta \vec{v}, \qquad
    \nabla \cdot \vec{v} = 0.
\]
L'existence et la r\'egularit\'e des solutions en dimension 3 reste l'un des
probl\`emes du mill\'enaire.

\paragraph{M\'ecanique quantique.}
L'\'equation de Schr\"odinger d\'ecrit l'\'evolution d'une particule quantique~:
\[
    i\hbar \frac{\partial \psi}{\partial t} = -\frac{\hbar^2}{2m}\Delta \psi + V\psi.
\]

Ces exemples illustrent la richesse et l'universalit\'e des EDP. Ce cours
fournira les outils math\'ematiques n\'ecessaires pour les comprendre et les
r\'esoudre.

\section*{R\'ef\'erences bibliographiques}

Ce cours s'inspire de plusieurs ouvrages de r\'ef\'erence~:
\begin{itemize}
    \item \textsc{Evans}, L.C., \emph{Partial Differential Equations}, AMS, 2010.
    \item \textsc{Br\'ezis}, H., \emph{Analyse fonctionnelle}, Dunod, 2005.
    \item \textsc{Zuily}, C. et \textsc{Queff\'elec}, H., \emph{Analyse pour
          l'agr\'egation}, Dunod.
    \item \textsc{Raviart}, P.-A. et \textsc{Thomas}, J.-M., \emph{Introduction
          \`a l'analyse num\'erique des EDP}, Dunod, 1998.
    \item \textsc{Strauss}, W., \emph{Partial Differential Equations: An Introduction},
          Wiley, 2008.
    \item \textsc{John}, F., \emph{Partial Differential Equations}, Springer, 1982.
\end{itemize}

\vspace{1cm}
\noindent
Bonne lecture et bon travail~!

\vspace{0.5cm}
\hfill\emph{L'auteur}
